\documentstyle{article}
\begin{document}

\centerline{\bf Quasi-exactly Solvable Spectral Problems for Ordinary
Differential Operators}

\vskip 0.3cm

\centerline{Niky Kamran}
\centerline{Department of Mathematics and Statistics}
\centerline{McGill University}
\centerline{Montreal, Quebec H3A 2K6, Canada}
\centerline{Email: nkamran@math.mcgill.ca}

\vskip 0.3cm

An ordinary differential operator is said to be quasi-exactly solvable if
it is expressible as a polynomial in the generators of a
finite-dimensional Lie algebra of first-order differential operators
admitting a finite-dimensional module of smooth functions. This Lie
algebra can thus be thought of as a hidden symmetry algebra for the
given operator. Indeed, if the module is $n$-dimensional, then one can
obtain $n$ eigenvalues of the operator, counting multiplicities, by
diagonalizing an $n \times n$ matrix, which is an algebraic problem. Many
Schr\"odinger operators of interest in quantum mechanics are quasi-exactly
solvable. Examples include some families of sextic anharmonic oscillators,
as well as the Morse, P\"oschl-Teller and Mathieu potentials. We
will survey the known results on the classification and normal forms for
quasi-exactly solvable ordinary differential operators and discuss some
open problems. This is joint work with Artemio Gonzalez-Lopez (Universidad
Complutense, Madrid) and Peter Olver (University of Minnesota).

\end{document}
